\subsection{闭子群}

定理 2. 设 G 是 R 或 \(\mathbf{Q}_{p}\)\(e\) 上的有限维李群。G 的每个闭子群都是 G 的李子群。更一般地，设 U 是 G 中 e 的对称开邻域，H 是 U 的非空闭子空间，且 \(x \in \mathrm{H}\)、\(y \in \mathrm{H}\) 和 \(xy^{-1} \in \mathrm{U}\) 蕴含 \(xy^{-1} \in \mathrm{H}\)。则 H 是 G 的李子群芽。

令 \(\mathfrak{h}\)\(e\) 为 H 在 e 处的切李子代数（§ 4，第 5 号，定义 2）。存在 G 的一个李子群芽 \(\mathrm{H}_0\)，其李代数为 \(\mathfrak{h}\)，且包含于 H。我们证明 \(\mathrm{H}_0\) 在由 G 诱导的拓扑下是 H 的开子集。这将证明 H 是 G 的解析子流形，并完成定理的证明。

存在与 \(\mathfrak{l}\) 在 \(\mathfrak{h}\) 中互补的向量子空间 \(\mathrm{L}(\mathrm{G})\)，在 \(\mathrm{V}_1, \mathrm{V}_2\)\(\mathfrak{h}\) 和 \(\mathfrak{l}\) 中分别有包含 0 的对称开邻域 ，以及定义在 \(\phi\) 上的 G 的指数映射 \(\mathrm{G}\)\(\mathrm{V}_1 + \mathrm{V}_2\)，具有如下性质：

\begin{enumerate}
\def\labelenumi{(\alph{enumi})}
\item
  映射 \((a_{1}, a_{2}) \mapsto \phi(a_{1})\phi(a_{2})\) 是从 \(\mathbf{V}_{1} \times \mathbf{V}_{2}\) 到 \(\mathbf{V}\) 的开子集 \(\mathbf{G}\) 的解析同构；
\item
  \(\phi (\mathrm{V}_1)\subset \mathrm{H}_0;\)；
\item
  \(V^{2} \subset U.\)。
\end{enumerate}

我们将证明（这也将完成证明），存在 \(\mathbf{V}_2'\)\(\mathbf{V}_2\) 中 0 的一个开邻域，并且使得 \(\mathbf{H} \cap (\phi(\mathbf{V}_1)\phi(\mathbf{V}_2')) = \phi(\mathbf{V}_1)\)。

假设断言不成立。于是可以找到 \((x_{n})\) 中的序列 \(\mathbf{V}_1\) 和 \((y_{n})\) 中趋于 0 的序列 \(\mathbf{V}_2 - \{0\}\)，使得对所有 n 都有 \(\phi (x_n)\phi (y_n)\in \mathbf{H}\)。由 (c)，则 \(n\)\(\phi (y_n)\in \mathbf{H}\)。

若 \(\mathrm{K} = \mathbf{Q}_p\)，可假设 \(\mathrm{V}_2\) 是 \(\mathfrak{l}\) 的加法子群，并且对所有 \(\phi (pa) = \phi (a)^p\) 和所有 \(a\in \mathrm{V}_2\)，有 \(p\in \mathbf{Z}\)。于是对所有 \(\phi (\lambda y_1)\in \mathbf{H}\)，都有 \(\lambda \in \mathbf{Z}\)，由连续性，对所有 \(\lambda \in \mathbf{Z}_p\) 亦然。从 \(f\colon \lambda \mapsto \phi (\lambda y_1)\) 到 G 的映射 \(\mathbf{Z}_p\) 是解析的，且取值于 H，并且 \(G\)\((T_0f)(1) = y_1\)。因此 \(y_{1}\in \mathfrak{h}\)，矛盾。于是定理在 \(\mathbf{Q}_p\) 的情形成立。

若 \(\mathbf{K} = \mathbf{R}\)，可假设 \(\mathrm{V}_2\) 连通，且 \(y_{n}\) 属于 \(\frac{1}{4}\mathrm{V}_2 - \{0\}\)。必要时取 \((y_n)\) 的子序列，可以找到非零标量序列 \((\lambda_n)\)，使得 \(\lambda_n^{-1}y_n\) 趋于 \(y\) 中的一个元素 \(\mathrm{V}_2 - \{0\}\)。序列 \((\lambda_n)\) 趋于 0。令 \(\lambda \in \mathbf{R}\) 满足 \(\lambda y \in \frac{1}{4}\mathrm{V}_2\)；我们证明 \(\exp (\lambda y) \in H\)。可以假设对所有 n 都有 \(\lambda\lambda_n^{-1}y_n \in \frac{1}{4}\mathrm{V}_2\)。令 \(n\)\(k_n \in \mathbf{Z}\)，使得 \(|\lambda - k_n\lambda_n|\) 趋于 0。当 n 充分大时，\(n\)\((\lambda - k_n\lambda_n)\lambda_n^{-1}y_n \in \frac{1}{4}\mathrm{V}_2\)，从而 \(k_n y_n \in \frac{1}{4}\mathrm{V}_2\)。因此，当 h 为整数且 \(\exp (hy_n) \in H\) 时，\(h\)\(0 \leqslant |h| \leqslant |k_n|\)（按 \(|h|\) 归纳即可看出）。于是

\[
\begin{array}{r l} \exp (\lambda y) & = \lim _ {n \to \infty} \exp (\lambda \lambda_ {n} ^ {- 1} y _ {n}) = \lim _ {n \to \infty} (\exp ((\lambda - k _ {n} \lambda_ {n}) \lambda_ {n} ^ {- 1} y _ {n}) \exp (k _ {n} y _ {n})) \\ & = \lim _ {n \to \infty} \exp k _ {n} y _ {n} \in H. \end{array}
\]

因此，映射 \(f \colon \lambda \mapsto \exp \lambda y\)（其中 \(\lambda y \in \frac{1}{4} \mathrm{V}_2\)）取值于 \(\mathbf{H}\)，且 \((\mathrm{T}_0f)(1) = y\)。所以 \(y \in \mathfrak{h}\)，矛盾。于是定理在 \(\mathbf{R}\) 的情形成立。

若不假设 G 有限维\(G\)，定理 2 不再成立（练习 12）。

 推论 1. 设  是局部紧群，G 是 R（分别地，\(\mathbf{Q}_{p}\)）上的有限维李群，f 是从  到 G 的连续态射。若 f 的核是离散的，则  是有限维实（分别地，p 进）李群。

存在  中 e 的一个紧邻域 V，使得 f\textbar V 是从 V 到 G 的一个紧子空间的同胚。若 U 是 G 中充分小的 e 的开邻域，则以 H = f(V) ∩ U 代入时满足定理 2 的假设。因此 H 是 G 的李子群芽。令 W 为 f\textbar{} V 下 H 的逆像。则 W 是  中 e 的一个邻域。用 f\textbar{} W 的逆映射将 H 上的解析流形结构搬到 W 上。对于所有 \(^{-1}\)\(z \in G'\)，从 G 到 G 的映射 \(x \mapsto f(z)xf(z)^{-1}\) 是解析的；因此 W 中存在 e 的一个开邻域 W'，使得从 W' 到 W 的映射 \(x' \mapsto zx'z^{-1}\) 是解析的。由 § 1，第 9 号，命题 18，G' 上存在一个李群结构，使得它在 e 的充分小开邻域上诱导出与 W 相同的解析结构，因而也与 G' 的初始拓扑相同。

推论 2. 设 G 是 K 上的有限维李群，H 是 G 的子群，V 是 G 中 e 的一个开邻域，\((\mathbf{M}_{i})_{i\in I}\) 是 K 上的一族解析流形；对于所有 \(i\in I\)，设 \(f_{i}\) 是从 V 到 \(M_{i}\) 的 K-解析映射，满足

\[
\mathrm{H} \cap \mathrm{V} = \{x \in \mathrm{V} | f _ {i} (x) = f _ {i} (e) \text {   for   all   } i \in \mathrm{I} \}.
\]

\begin{enumerate}
\def\labelenumi{(\roman{enumi})}
\item
  若 \(\mathbf{K} = \mathbf{C}\)，则 \(\mathbf{H}\) 是 \(\mathbf{G}\) 的李子群。
\item
  若 \(\mathbf{K}\) 是 \(\mathbf{Q}_p\) 的有限扩张且 \(\mathbf{I}\) 有限，则 \(\mathbf{H}\) 是 \(\mathbf{G}\) 的李子群。
\item
  假设 K = C。将 G 看作实李群。则 H 是 G 的实李子群（定理 2）。令 \(a \in \mathbf{L}(\mathbf{H})\)。存在 C 中 0 的连通开邻域 W，使得对所有 \(\exp \lambda a \in V\)，\(\lambda \in W\)。令 \(i \in I\)。若 \(f_i(\exp \lambda a) = f_i(e)\)，则 \(\lambda \in R \cap W\)。由解析延拓，对所有 \(f_i(\exp \lambda a) = f_i(e)\)，都有 \(\lambda \in W\)。因此当 \(\exp \lambda a \in H\) 时 \(\lambda \in W\)，从而对所有 \(\mu a \in L(H)\)，都有 \(\mu \in C\)。所以 H 是复李群 G 的李子群（§ 4，第 2 号，命题 2）。
\item
  假设 K 是 \(Q_{p}\)\(Q_{p}\) 的有限扩张。将 G 看作  上的李群。它是有限维的，且由定理 2，H 是 G 的 p 进李子群。由于 I 有限，\(\prod_{i\in I} M_{i}\) 是一个流形，可以假设族 \((f_{i})\) 归结为一个映射 f。~令 \(a \in L(G)\)。令 \(\phi\) 为 G 的一个指数映射。于是当 \(f(\phi(\lambda a)) = f(e)\) 且 \(\lambda \in Q_{p}\) 充分小时，\(|\lambda|\)。由于 f 是 K-解析的，可知当 \(f\)\(f(\phi (\lambda a)) = f(e)\) 且 \(\lambda \in \mathbf{K}\) 充分小时，\(|\lambda|\)。因此当  且 \(\phi (\lambda a) \in \mathrm{H}\) 充分小时，\(\lambda \in \mathrm{K}\)\(|\lambda|\)，从而对所有 ，都有 \(\mu a \in \mathrm{L(H)}\)\(\mu \in \mathrm{K}\)。证明如 (i) 完成。
\end{enumerate}

推论 2 (ii) 在省去 I 有限这一假设时不再成立。

